\section{Introduction}

% provenance: dec_01KXT240A70R8N104BPN5G07CJ (RQ1) motivates the reach/scope question.
Distributed energy resources now participate in grid operations at a scale that makes their
control interface a systemic concern. IEEE~2030.5 (Smart Energy Profile~2.0), profiled for
distribution use by the Common Smart Inverter Profile (CSIP), is the application protocol through
which utilities and aggregators dispatch inverter-based resources. An adversary who obtains a
credentialed identity in this ecosystem does not merely read data; it commands real and reactive
power on physical hardware connected to the grid.

% provenance: jrn_01KXT3JPAZ8RSGZK99X4G3HWPE (SAND2019-1490 4.1) supports the two facts here.
Two properties of the trust model make this harder than a conventional credential-theft problem.
The device's bootstrap identity and its operational authorization are the same certificate, so
there is no separate short-lived credential to lose. And that certificate does not expire and
cannot be revoked: the analysis of record for these standards states that the PKI is distinguished
by ``the use of non-expiring certificates and the explicit prohibition of CRL and OCSP'', and that
``if the device's private key is compromised, the device certificate cannot be
revoked''~\cite{obert2019recommendations}. What the standard offers instead is a local allow or
deny list.

The obvious objection is that a denylist is enough --- the operator refuses the compromised
identity and the incident is over. \textbf{We measure that assumption and it does not hold.}
Denying the identity of an adversary who is already connected reduces integrated feeder harm by
\emph{exactly zero} on both feeders we evaluate. It refuses new sessions, which is all it claims
to do; the adversary's established session continues to issue commands, and the scheduled control
it has already installed continues to act on the inverter. Containment fails not because authority
cannot be withdrawn but because withdrawing it at one layer leaves two others running.

That observation organises the paper. Bounding a DER compromise requires controlling two
independent quantities: the \emph{set of physical actions} an authorization admits, and the
\emph{duration} for which authority remains effective at every layer that can carry it --- the
credential, the open session, and the issued command.

CONTAINDER separates bootstrap identity from short-lived, attested, scope-bound operational
authority, and enforces expiry at all three layers. Because grid actuation persists in a way
computation does not, expiry that stops at the credential is not containment: we show that a
credential lifetime with neither session nor command enforcement removes no harm at all, and that
command cleanup without session enforcement removes none either, because the adversary re-issues
over the session that is still open.

\subsection*{Contributions}
\begin{itemize}
\item \emph{A DER credential and enforcement architecture} that separates bootstrap identity,
operational credential and value-constrained authorization, and enforces expiry at the credential,
session and command-effect layers, deployable through a gateway proxy that leaves downstream
IEEE~2030.5 endpoints unmodified.

\item \emph{A measurement of what each containment layer is worth.} Identity denial alone removes
$0\%$ of integrated harm; adding session termination removes $48.8$--$54.4\%$; adding command
cancellation removes $70.6$--$76.3\%$. The two enforcement mechanisms fail alone for different
reasons and are separately necessary --- a claim an earlier version of this work asserted from a
harness that could not identify it, and which is here earned from a lifecycle model that
distinguishes the three layers.

\item \emph{An authorization-shape result, with the condition under which it matters.} At
identical width, a reactive bound that admits zero absorption fails to contain where an absorption
floor contains, at negligible cost to volt-var service. The magnitude of that effect scales with
the reactive absorption the feeder is actually drawing rather than with PV penetration; it
survives at an operating point where legitimate operation is compliant, but is an order of
magnitude smaller there than at an overstressed one.

\item \emph{A corrected, released evaluation harness} on two public feeders, with a
pre-registration committed to version control before any treatment arm ran, every departure from
it disclosed, and two defects in the predecessor harness identified and repaired.
\end{itemize}

We claim no novelty for short-lived certificates or for expiry-based containment. The Web PKI
ratified that trade in 2023~\cite{cabf-sc063v4}, and the Sandia analysis of these very standards
recommends limited-life certificates and records that manufacturers are ``not required to limit
the life of a device certificate''~\cite{obert2019recommendations} --- which makes a short-lived
operational credential permitted rather than new. What is new here is the mapping onto DER
authorization: which layers expiry must reach before it contains anything, which control point an
authorization must withhold, and the operating condition that decides whether either matters.
